\documentclass[draft]{agujournal2019}

\usepackage[T1]{fontenc}
\usepackage[utf8]{inputenc}
\usepackage{textcomp}
\usepackage{url}
\usepackage{lineno}
\usepackage{amsmath,amssymb}
\usepackage{booktabs}
\usepackage{longtable}
\usepackage{array}
\usepackage{graphicx}
\makeatletter
\define@key{Gin}{alt}{}
\makeatother
\usepackage{hyperref}

\providecommand{\tightlist}{%
  \setlength{\itemsep}{0pt}\setlength{\parskip}{0pt}}
\providecommand{\pandocbounded}[1]{#1}

\linenumbers
\draftfalse
\journalname{JGR: Solid Earth}

\begin{document}

\title{Seismic Velocity Changes as Stress and Strain Meters: A Unified
Framework for Environmental, Tectonic, and Volcanic Monitoring}

\authors{Marine A. Denolle\affil{1}}

\affiliation{1}{Department of Earth and Space Sciences, University of Washington, Seattle, WA, USA}

\correspondingauthor{Marine A. Denolle}{mdenolle@uw.edu}

\begin{keypoints}
\item A unified framework links seismic velocity changes to stress
strain hydrology and rheology
\item Directional crack fabric explains Parkfield velocity changes
tracking axial contraction
\item Coda window metrics make velocity change processing reproducible
for codameter
\end{keypoints}

\begin{abstract}
Ambient seismic noise monitoring of relative seismic velocity changes
(\(\delta v/v\)) has become a powerful tool for probing the subsurface
mechanical state across environmental, tectonic, and volcanic settings.
Yet interpretation remains fragmented: thermoelastic, hydrological,
poroelastic, and tectonic models are applied in isolation, making it
difficult to separate competing signals or jointly invert for subsurface
rheology. Here we develop a unified theoretical framework that treats
\(\delta v/v\) simultaneously as a stress meter and a strain meter,
grounded in nonlinear (third-order) elasticity, poroelastic theory, and
the induced-stress formulation of Tromp and Trampert (2018). Through
systematic scenario modeling we quantify the sensitivity of
\(\delta v/v\) to surface temperature variations
(\textasciitilde{}0.01--0.3\% annually), hydrological loading and
pore-pressure diffusion (\textasciitilde{}0.01--0.1\%), ice and snow
loading (\textasciitilde{}0.001--0.5\%), partially saturated media with
dynamic capillary effects, and tectonic strain accumulation
(\textasciitilde{}0.001--0.01\%/yr). We show that stress-induced
anisotropy from oriented microcrack closure explains why \(\delta v/v\)
correlates with contractional axial strain rather than dilatation, as
observed at Parkfield. We demonstrate the framework quantitatively at
three contrasting sites: at Parkfield (strike-slip), the isotropic
formulation fails and the deviatoric framework predicts a stress
accumulation rate of \textasciitilde12 kPa/yr at \textasciitilde0.8 km
depth in fractured granite (\(\beta \approx 240\),
\(\mu' \approx 250\)); at Cascadia (subduction), the isotropic framework
succeeds and, using the borehole strain calibration, gives 0.58 kPa/yr
in marine sediments (\(\beta \approx 3160\), \(\mu' \approx 620\)); and
at Kīlauea (volcanic caldera collapse), concentric ring fractures cause
\(\delta v/v\) to track radial rather than volumetric strain, with
\(\beta_{\text{radial}} \approx 300\) giving a reservoir-stress estimate
within a factor of two of geodetic constraints. Because \(\delta v/v\)
depends on both frequency and coda lapse time --- with different
frequencies and windows sampling different wave types and depth ranges
--- multi-frequency, multi-lapse \(\delta v/v\) combined with geodetic
strain from GNSS/InSAR has the potential to enable depth-resolved stress
and strain imaging, but only with appropriate kernels, regularization,
and independent constraints. We map the parameter space where key
assumptions remain valid, identify diagnostic signatures that
distinguish elastic, viscoelastic, and slow-dynamics rheologies, and
propose a reproducible window-selection framework for constraining
subsurface rheological models. Companion notebooks reproduce the main
and supporting figures and provide practical computational tools for
implementing this framework.
\end{abstract}

\section*{Plain Language Summary}
Small changes in seismic wave speed can reveal how rocks, sediments, and
fluids respond to environmental and tectonic forcing. This paper
connects several explanations that are often treated separately,
including temperature, groundwater, surface loading, earthquake damage,
volcanic pressure changes, and long term tectonic strain. The result is
a practical framework for deciding when velocity changes behave like a
stress meter, when they behave like a strain meter, and when additional
material or hydrological effects must be modeled.

\section{1. Introduction}\label{introduction}

\subsection{1.1 The Promise of Ambient Noise
Monitoring}\label{the-promise-of-ambient-noise-monitoring}

The continuous monitoring of relative seismic velocity changes
(\(\delta v/v\)) using ambient seismic noise has transformed our ability
to observe the evolving mechanical state of the Earth's crust. Since
Sens-Schönfelder and Wegler (2006) first detected seasonal velocity
changes at Mt. Merapi correlated with groundwater fluctuations, the
technique has been applied across a remarkable breadth of settings:
active faults (Brenguier et al., 2008a; Taira et al., 2015; Okubo et
al., 2024), volcanoes (Brenguier et al., 2008b; Hotovec-Ellis et al.,
2022), groundwater systems (Clements and Denolle, 2018, 2023; Mao et
al., 2022), urban basins (Ermert et al., 2023), glaciated regions
(Donaldson et al., 2019), arid deserts (Richter et al., 2014), and
agricultural soils (Shi et al., 2026).

The physical basis is that coda waves --- scattered waves arriving after
ballistic arrivals --- are exquisitely sensitive to small perturbations
in the medium (Lobkis and Weaver, 2003; Snieder et al., 2002). Ambient
noise cross-correlation (Nakata et al., 2019) removes the dependence on
earthquake sources, enabling continuous monitoring from seconds to
decades. The quantity \(\delta v/v\) provides a relative measure of the
volume-averaged perturbation in seismic velocity, with depth sensitivity
controlled by the frequency content of the measurement (Obermann et al.,
2014).

\subsection{1.2 The Interpretation
Challenge}\label{the-interpretation-challenge}

The observed \(\delta v/v\) time series at any station is a
superposition of signals from thermoelastic stresses (Richter et al.,
2014; Lecocq et al., 2017), hydrological effects including pore-pressure
changes (Clements and Denolle, 2018, 2023; Fokker et al., 2021), surface
loading from water, ice, and snow (Tsai, 2011; Donaldson et al., 2019),
capillary effects in partially saturated media (Shi et al., 2026),
coseismic damage and post-seismic healing (Wegler and Sens-Schönfelder,
2007; Brenguier et al., 2008a), and tectonic strain accumulation (Okubo
et al., 2024).

The 20-year \(\delta v/v\) record at Parkfield, California (Okubo et
al., 2024) exemplifies the challenge: the time series contains seasonal
oscillations from environmental effects, coseismic drops from the 2003
San Simeon and 2004 Parkfield earthquakes, logarithmic post-seismic
healing, and a statistically significant long-term increase correlating
with inter-seismic tectonic strain. In Oklahoma, Zhang et al.~(2023)
showed that multi-year hydrological effects dominate the overall
\(\delta v/v\), while annual thermoelastic strains dominate the seasonal
cycle. In Greece, Delouche et al.~(2023) identified hydrological cycling
as the primary driver of seasonal \(\delta v/v\) at depth, ruling out
thermoelastic strain based on the frequency dependence of the signal.

\subsection{1.3 What Has Been Missing}\label{what-has-been-missing}

Several gaps persist:

\begin{enumerate}
\def\labelenumi{\arabic{enumi}.}
\item
  \textbf{Fragmented theory.} The thermoelastic model (Berger, 1975;
  Richter et al., 2014), poroelastic model (Roeloffs, 1988; Clements and
  Denolle, 2023), and induced-stress framework (Tromp and Trampert,
  2018) have been developed independently. A unified treatment through
  nonlinear elasticity is lacking.
\item
  \textbf{Systematic sensitivity analysis.} While individual studies fit
  models to specific datasets, a comprehensive parameter-space
  exploration has not been presented.
\item
  \textbf{Partially saturated media.} Most models assume fully saturated
  or fully drained conditions, yet Shi et al.~(2026) showed that dynamic
  capillary effects in partially saturated soils produce distinct
  hysteretic \(\delta v/v\) signatures not captured by existing
  frameworks.
\item
  \textbf{Anisotropy.} Most interpretations assume isotropic media, yet
  the Parkfield observations (Okubo et al., 2024) and the Tromp and
  Trampert (2018) theory show that deviatoric stress creates
  directionally dependent velocity changes.
\item
  \textbf{Depth-resolved inversion.} The frequency and lapse-time
  dependence of \(\delta v/v\) --- with different frequencies and coda
  windows sampling different wave types and depths --- represents an
  underexploited opportunity for 3-D stress/strain imaging when combined
  with geodetic data.
\item
  \textbf{Processing reproducibility.} There is little field-wide
  consensus on frequency bands, substack length, and the start and end
  of the coda window. These choices are often treated as processing
  details even though they change the sensitivity kernel and therefore
  the physical interpretation.
\item
  \textbf{Rheological inversion.} A systematic framework for jointly
  inverting \(\delta v/v\) and geodetic strain for subsurface rheology
  has not been formulated.
\end{enumerate}

\subsection{1.4 Scope and Contribution}\label{scope-and-contribution}

This paper addresses these gaps by developing a unified framework
connecting \(\delta v/v\) to stress, strain, and rheology through
nonlinear elasticity. Figure 1 summarizes the workflow, from forcing and
sensitivity kernels to mechanism diagnosis and the future
\texttt{codameter} implementation package. We present the theory
(Section 2), quantitative scenario models for thermoelastic (Section 3),
hydrological including partially saturated (Section 4), nonlinear
elastic (Section 5), and anisotropic (Section 6) effects. We then
develop the rheological inversion framework with emphasis on
depth-resolved imaging and reproducible coda-window selection (Section
7), map validity ranges (Section 8), apply the framework quantitatively
to three contrasting sites --- Parkfield (strike-slip), Cascadia
(subduction), and Kīlauea (volcanic) --- predicting stress at depth from
published \(\delta v/v\) measurements (Section 9), and discuss
implications (Section 10).

\begin{center}\rule{0.5\linewidth}{0.5pt}\end{center}

\section{\texorpdfstring{2. Unified Theoretical Framework:
\(\delta v/v\) as Both Stress and Strain
Meter}{2. Unified Theoretical Framework: \textbackslash delta v/v as Both Stress and Strain Meter}}\label{unified-theoretical-framework-delta-vv-as-both-stress-and-strain-meter}

\subsection{\texorpdfstring{2.0 Notation: What Does \(\delta v/v\)
Measure?}{2.0 Notation: What Does \textbackslash delta v/v Measure?}}\label{notation-what-does-delta-vv-measure}

The quantity \(\delta v/v\) (equivalently written \(dv/v\) in much of
the literature) is the relative velocity change measured from coda-wave
interferometry (CWI) via the relation \(\delta t/t = -\delta v/v\)
(Poupinet et al., 1984; Clarke et al., 2011; Snieder et al., 2002). A
critical question is: \emph{what velocity does \(\delta v/v\)
represent?}

Snieder (2002) showed that in a multiply scattering elastic medium,
P-to-S mode conversions drive the wavefield toward equipartition, where
the energy ratio \(E_S/E_P = 2(V_P/V_S)^3\). For a Poisson solid
(\(V_P/V_S = \sqrt{3}\)), S-waves carry \textasciitilde{}90\% of the
scattered energy. Consequently, the CWI measurement is dominated by
changes in S-wave velocity: \(\delta v/v \approx \delta V_S/V_S\). Singh
et al.~(2019) generalized this, showing that the CWI velocity change
varies along the coda as a weighted sum
\([\delta V/V]_{\text{CWI}} = q(t,\gamma)\,\delta V_S/V_S + [1-q(t,\gamma)]\,\delta V_P/V_P\),
where the S-wave weight \(q \to 1\) at late coda times.

When surface waves dominate the ambient noise cross-correlations --- as
is typical at periods \textgreater{} 1 s --- the measured \(\delta v/v\)
corresponds to a phase-velocity change \(\delta c/c\) (Ermert et al.,
2023; Fokker et al., 2021). This relates to the depth profile of
\(\delta V_S / V_S\) through the sensitivity kernel (Fokker et al.,
2021, their Eq. 13):

\[\frac{\delta v}{v}(f, t) \approx \frac{\delta c}{c}(f, t) = \frac{1}{c(f)} \int_0^\infty \frac{\partial c}{\partial V_S}(z, f) \, \delta V_S(z, t) \, dz \tag{1}\]

Throughout this paper we adopt the following notation:

{\def\LTcaptype{none} % do not increment counter
\begin{longtable}[]{@{}
  >{\raggedright\arraybackslash}p{(\linewidth - 2\tabcolsep) * \real{0.4706}}
  >{\raggedright\arraybackslash}p{(\linewidth - 2\tabcolsep) * \real{0.5294}}@{}}
\toprule\noalign{}
\begin{minipage}[b]{\linewidth}\raggedright
Symbol
\end{minipage} & \begin{minipage}[b]{\linewidth}\raggedright
Meaning
\end{minipage} \\
\midrule\noalign{}
\endhead
\bottomrule\noalign{}
\endlastfoot
\(\delta v / v\) & Measured relative velocity change (the observable) \\
\(\delta V_S / V_S\) & Local relative change in shear-wave velocity at
depth \(z\) (inferred) \\
\(\delta c / c\) & Relative change in surface-wave phase velocity at
frequency \(f\) \\
\end{longtable}
}

The Fokker et al.~(2021) notation \(d\beta/\beta\) refers to
\(\delta V_S/V_S\); we use the latter for clarity. Note that Equation 1
establishes that \(\delta v/v\) is a \emph{depth-averaged},
\emph{frequency-dependent} measurement --- different frequencies probe
different depths, a property we exploit for 3-D imaging (Section 7).

\subsection{2.1 Nonlinear Elasticity: The Murnaghan
Framework}\label{nonlinear-elasticity-the-murnaghan-framework}

The sensitivity of seismic velocities to stress originates in nonlinear
elastic behavior. Murnaghan (1937) showed that expanding the elastic
energy density to third order in strain introduces three additional
constants (\(l\), \(m\), \(n\)):

\[\rho_0 \phi = \frac{\lambda + 2\mu}{2} I_1^2 - 2\mu I_2 + l I_1^3 + m I_1 I_2 + n I_3 \tag{2}\]

These third-order elastic (TOE) constants cause seismic velocities to
depend on the ambient stress state. The acoustoelastic relation
(Ostrovsky and Johnson, 2001; Clements and Denolle, 2023) gives:

\[\frac{\delta v}{v} \approx \beta \, \epsilon_{kk}, \quad \beta = \frac{3}{2} + \frac{l + 2m}{\lambda + 2\mu} \tag{3}\]

where \(\beta\) ranges from \textasciitilde{}−10 (steel) to
\textasciitilde{}−10\(^4\) (sandstone), reflecting microstructural
compliance (Hughes and Kelly, 1953; Johnson and Rasolofosaon, 1996;
Clements and Denolle, 2023).

\subsection{2.2 The Induced-Stress
Formulation}\label{the-induced-stress-formulation}

Tromp and Trampert (2018) showed that induced stress modifies shear-wave
velocity as:

\[\frac{\delta V_S}{V_S} = \frac{\mu' p^0}{2\mu} + \frac{1 - \mu'}{4\mu} \hat{\mathbf{k}} \cdot \boldsymbol{\tau}^0 \cdot \hat{\mathbf{k}} - \frac{1 + \mu'}{4\mu} \hat{\mathbf{a}} \cdot \boldsymbol{\tau}^0 \cdot \hat{\mathbf{a}} \tag{4}\]

where \(\hat{\mathbf{k}}\) and \(\hat{\mathbf{a}}\) are propagation and
polarization directions, \(p^0\) is induced pressure,
\(\boldsymbol{\tau}^0\) is deviatoric stress, and \(\mu' = d\mu/dP\).
Critically, \emph{deviatoric stress creates anisotropy} without
requiring third-order elasticity.

\subsection{2.3 Poroelastic Extension}\label{poroelastic-extension}

Fokker et al.~(2021) combined Equation 4 with poroelastic theory for
surface loading \(T_{33}^0\) and pore pressure \(u^0\). For horizontally
propagating surface-wave endmembers under a positive-compressive load
convention, their coefficients reduce to:

\[\left(\frac{\delta V_S}{V_S}\right)_{\mathrm{Rayleigh/SV}} = -\frac{\mu'}{2\mu}u^0 + \frac{\mu' + 1}{12\mu}T_{33}^0,\qquad
\left(\frac{\delta V_S}{V_S}\right)_{\mathrm{Love/SH}} = -\frac{\mu'}{2\mu}u^0 . \tag{5}\]

The pore-pressure term decreases velocity by reducing effective stress.
The vertical-load term depends on propagation and polarization geometry;
therefore component-specific applications should evaluate the full
directional form in Equation 4 rather than using a scalar loading
coefficient.

\subsection{\texorpdfstring{2.4 General \(\delta v/v\)
Model}{2.4 General \textbackslash delta v/v Model}}\label{general-delta-vv-model}

The observed \(\delta v/v\) is modeled as (Okubo et al., 2024):

\[\frac{\delta v}{v}(t) = s_T \cdot T(t - t_{\text{shift}}) + p_1 \cdot \Delta\text{GWL}(t) + \sum_i s_i L_i(t) + b_0 t + a_0 \tag{6}\]

where each term has a clear physical origin: thermoelastic strain
(Richter et al., 2014), hydrological pore-pressure change
(Sens-Schönfelder and Wegler, 2006), coseismic damage with logarithmic
healing (Snieder et al., 2017), and long-term tectonic loading. The
unifying insight is that all terms operate through stress-dependent
modification of elastic moduli.

\subsection{2.5 Bridging the Strain and Stress
Formulations}\label{bridging-the-strain-and-stress-formulations}

A central claim of this paper is that the strain formulation (Eq. 3,
from acoustoelasticity) and the stress formulation (Eq. 4, from Tromp
and Trampert, 2018) describe the same physics from different starting
points. We can make this explicit. Under hydrostatic (isotropic)
conditions, the induced pressure relates to volumetric strain through
the bulk modulus: \(p^0 = -\kappa \, \epsilon_{kk}\). Substituting into
the isotropic term of Equation 4 gives
\(\delta V_S/V_S = -\mu' \kappa \, \epsilon_{kk} / (2\mu)\). Equating
with the acoustoelastic relation (Eq. 3,
\(\delta v/v = \beta \, \epsilon_{kk}\)) and using
\(\delta v/v \approx \delta V_S/V_S\) yields:

\[\beta = -\frac{\mu' \kappa}{2\mu} \tag{7}\]

This important relation connects the ``materials science'' parameter
\(\beta\) --- measured in laboratory acoustoelastic experiments via
third-order elastic constants (Hughes and Kelly, 1953) --- to the
``seismological'' parameter \(\mu'\) --- inferred from the depth profile
of shear modulus versus confining pressure (Fokker et al., 2021). For
typical shallow sediments (\(\mu' \approx 80\), \(\kappa \approx 5\)
GPa, \(\mu \approx 0.5\) GPa), Equation 7 gives \(\beta \approx -400\),
consistent with the range reported for weakly cemented rocks (Clements
and Denolle, 2023).

An important caveat is that Equation 7 uses the \emph{drained} bulk
modulus \(\kappa\), appropriate for loading timescales longer than the
pore-pressure diffusion time. At shorter timescales (the undrained
limit), the effective bulk modulus is
\(\kappa_u = \kappa / (1 - \alpha_B B)\), where \(\alpha_B\) is the Biot
coefficient and \(B\) is the Skempton coefficient (Roeloffs, 1988). In
this limit the effective \(\beta\) is larger in magnitude --- meaning
the velocity is \emph{more} sensitive to strain when the pore fluid has
not had time to drain. This drained-undrained distinction connects
directly to the poroelastic models in Section 4.1 and explains why the
\(\delta v/v\) response to a rapid loading event (e.g., rainfall) may
differ in amplitude from the response to a slow loading event (e.g.,
seasonal groundwater) even for the same total strain change.

Crucially, \emph{which} modulus applies is not an assumption but a
property of the data: it is set by the Péclet number
\(\mathrm{Pe} = \omega L^2/c\) comparing the forcing angular frequency
\(\omega\) to the drainage rate \(c/L^2\) at the coda sensitivity depth
\(L = V_S/(3f)\), for hydraulic diffusivity \(c\). \(\mathrm{Pe} \ll 1\)
is drained (use \(\kappa\)), \(\mathrm{Pe} \gg 1\) is undrained (use
\(\kappa_u\)), and \(\mathrm{Pe}\sim 1\) is transitional. Because the
seismic-band moduli inferred from \(V_P, V_S\) are themselves the
\emph{undrained} moduli (the pore fluid cannot drain within a
\textasciitilde{}1 Hz wave period), the natural workflow is to take
\(\kappa_u\) from velocities and, where the signal is drained, convert
via \(\kappa = \kappa_u(1-\alpha_B B)\). We apply this test explicitly
at each site in Section 9; for example, the Cascadia secular trend has
\(\mathrm{Pe}\sim 2.5\) (transitional-to-undrained), so \(\kappa_u\)
from the seismic velocities is the appropriate modulus there.

The equivalence holds \emph{only under isotropic loading}. Under
deviatoric stress, the Tromp and Trampert formulation (Eq. 4) captures
directional effects through the
\(\hat{\mathbf{k}} \cdot \boldsymbol{\tau}^0 \cdot \hat{\mathbf{k}}\)
and
\(\hat{\mathbf{a}} \cdot \boldsymbol{\tau}^0 \cdot \hat{\mathbf{a}}\)
terms, which have no counterpart in the scalar acoustoelastic relation.
This divergence is precisely what makes the Parkfield observation ---
where \(\delta v/v\) correlates with contractional but not dilatational
strain (Okubo et al., 2024) --- a diagnostic for deviatoric stress
contributions.

\subsection{2.6 Alternative and Complementary
Mechanisms}\label{alternative-and-complementary-mechanisms}

Beyond the nonlinear-elastic framework, several alternative mechanisms
can produce temporal velocity changes. \textbf{Density changes} from
compaction can modify \(V_S\) through
\(\delta V_S/V_S = -\frac{1}{2}\delta\rho/\rho\), though Fokker et
al.~(2021) showed this is 2--3 orders of magnitude smaller than the
shear-modulus mechanism for typical settings. \textbf{Fluid
substitution} --- changes in pore-fluid composition (e.g., CO\(_2\)
injection; Zhu et al., 2019) --- modifies the bulk modulus via
Gassmann's equation and primarily affects \(V_P\) rather than \(V_S\).
\textbf{Mineral alteration and cementation} (hydration, dissolution,
bio-cementation) permanently modify the elastic frame on timescales of
months to years (Rodríguez Tribaldos and Ajo-Franklin, 2021).
\textbf{Direct temperature effects} on elastic moduli --- distinct from
the thermoelastic \emph{stress} mechanism --- are usually small
(\textasciitilde{}0.01\%/K) but become significant near phase
transitions such as freezing/thawing in permafrost (James et al., 2017).
\textbf{Scatterer relocation} from fracture opening or fluid migration
produces apparent \(\delta v/v\) that reflects structural change rather
than bulk velocity change; Obermann et al.~(2013) developed theory
separating these contributions. Finally, \textbf{source-side effects}
--- seasonal changes in ocean-wave spectra or anthropogenic noise ---
can produce spurious \(\delta v/v\) if not mitigated (Zhan et al., 2013;
Okubo et al., 2024).

\begin{center}\rule{0.5\linewidth}{0.5pt}\end{center}

\section{3. Thermoelastic Effects}\label{thermoelastic-effects}

\subsection{3.1 Temperature Diffusion and Thermoelastic
Stress}\label{temperature-diffusion-and-thermoelastic-stress}

Surface temperature fluctuations diffuse into the Earth governed by the
heat equation (Berger, 1975):

\[T(z,t) = T_0 \, e^{-\gamma z} \cos(\omega t - \gamma z), \quad \gamma = \sqrt{\omega/(2\kappa_T)} \tag{8}\]

with thermal skin depths of 1.2--4.5 m annually for
\(\kappa_T = 0.15\)--\(2.0\) mm\(^2\)/s (Richter et al., 2014; Ermert et
al., 2023; Figure S1). Daily variations penetrate only 0.06--0.24 m,
affecting \(\delta v/v\) at frequencies above 2--4 Hz, whereas annual
variations dominate the seasonal signal at 0.5--2 Hz.

The thermoelastic stress induced by these temperature changes contains
two terms with fundamentally different depth behaviors (Berger, 1975;
Richter et al., 2014, Eq. 11; Figure S2). The first term decays with the
thermal skin depth \(1/\gamma\) and represents \emph{direct thermal
stressing}: heated rock expands against the confinement of surrounding
cooler rock, building compressive stress near the surface. The second
term decays with the much longer scale \(1/k\) (where \(k = 2\pi/L\) is
the horizontal wavenumber of the temperature field, \(L \sim 10\) km)
and represents the \emph{mechanical equilibrium response} --- a
broad-scale stress adjustment needed to satisfy the free-surface
boundary condition, penetrating to kilometer depths with a constant
phase delay of \(5\pi/4\) relative to surface temperature. For annual
variations, \(k/\gamma \approx 10^{-3}\), confirming that the two terms
operate at well-separated depth scales. Tsai (2011) showed that this
thermoelastic framework also explains seasonal GPS position changes,
establishing a direct link between geodetically observed strain and
seismically observed velocity change.

\subsection{3.2 From Stress to Velocity
Change}\label{from-stress-to-velocity-change}

The thermoelastic stress converts to velocity change via (Richter et
al., 2014, Eq. 12):

\[\frac{\delta v}{v} = b \frac{\partial(\rho v^2)}{\partial \sigma_c} \frac{(1-\nu)\Delta\sigma_c}{E} \tag{9}\]

where \(b \approx 1.5\) for S-waves at \(\nu = 0.2\) (Birch, 1961). The
combined temperature sensitivity (Ermert et al., 2023, their Eq. 3) is
\(s_T = 2b\alpha \, \partial(\rho v^2)/\partial\sigma_c\), where
\(\alpha\) is the thermal expansion coefficient.

\subsection{3.3 Sensitivity Analysis}\label{sensitivity-analysis}

Our sensitivity analysis (Figure S3) reveals that
\(\partial(\rho v^2)/\partial\sigma_c\) --- ranging from
\textasciitilde{}50 (intact rock) to \textasciitilde{}1000
(salt-cemented sediments at PATCX; Richter et al., 2014) --- is the
dominant control, varying surface \(\delta v/v\) by a factor of 20. The
Poisson's ratio dependence is monotonic (higher \(\nu\) increases the
confinement factor \(b\)), and the depth profile of thermoelastic
\(\delta v/v\) is exponential with an e-folding depth equal to the
thermal skin depth. Ermert et al.~(2023) expand the surface temperature
into five Fourier harmonics to capture sub-annual variations in Mexico
City (Figure S4). At Parkfield, Okubo et al.~(2024) constrain the
thermal diffusion delay \(t_{\text{shift}}\) to 0--90 days. In the San
Jacinto fault zone, Hillers et al.~(2015) found that thermoelastic
strain dominates the seasonal \(\delta v/v\) in this arid setting,
whereas in Oklahoma, Zhang et al.~(2023) showed that thermoelastic and
hydrological contributions are comparable.

\begin{center}\rule{0.5\linewidth}{0.5pt}\end{center}

\section{4. Hydrological Loading, Pore Pressure, and Capillary
Effects}\label{hydrological-loading-pore-pressure-and-capillary-effects}

\subsection{4.1 Pore-Pressure Diffusion in Saturated
Media}\label{pore-pressure-diffusion-in-saturated-media}

Roeloffs (1988) derived the poroelastic response at depth to a surface
load, containing undrained and drained terms that depend on hydraulic
diffusivity \(c\) (varying over five orders of magnitude), Skempton
coefficient \(B\), and undrained Poisson's ratio \(\nu_u\) (Clements and
Denolle, 2023, Eq. 8--9; Figure S5). Talwani et al.~(2007) extended this
to time-series precipitation loading.

\subsection{4.2 Groundwater Level
Models}\label{groundwater-level-models}

Okubo et al.~(2024) model \(\Delta\)GWL from precipitation using
exponential decay (their Eq. 4; Figure S6). Clements and Denolle (2023)
compared multiple hydrological models across California, finding the
drained poroelastic model explains 48\% of sites, with basin sites
showing lower diffusivities consistent with longer groundwater
retention. The CDM\(_k\) empirical model (Smail et al., 2019) provides a
useful parameter-free alternative.

\subsection{4.3 Competition Between Loading and Pore
Pressure}\label{competition-between-loading-and-pore-pressure}

Fokker et al.~(2021) showed that surface loading and pore-pressure
diffusion produce \emph{opposing} \(\delta v/v\) effects. Our scenarios
(Figure 3) quantify this for ice loading (loading dominates for
impermeable media), seasonal rainfall (pore pressure typically
dominates, consistent with anti-correlation observed by Clements and
Denolle, 2023), and reservoir impoundment (balance evolves with drainage
time).

\subsection{4.4 Partially Saturated Media and Dynamic Capillary
Effects}\label{partially-saturated-media-and-dynamic-capillary-effects}

The frameworks of Sections 4.1--4.3 assume fully saturated conditions
below the water table. However, much of the near-surface is partially
saturated, and the mechanics of the vadose zone differ fundamentally
from the saturated case. Shi et al.~(2026) recently demonstrated that
\emph{dynamic capillary effects} govern the transient stiffness response
of partially saturated soils, producing hysteretic \(\delta v/v\)
signatures not captured by standard poroelastic models.

In partially saturated granular media, the shear-wave velocity depends
on the effective shear modulus \(\mu_{\text{eff}}\), which Shi et
al.~(2026) compute from contact mechanics:

\[v_S = \sqrt{\mu_{\text{eff}} / \rho} \tag{10}\]

with \(\mu_{\text{eff}}\) controlled by an effective confining pressure
that includes both overburden and capillary stress. A compact way to
express the dynamic-capillary contribution is:

\[P_e(z,t) = \sigma_g'(z,t) + P_c^{\mathrm{eq}}(S_w) + \tau(S_w)\frac{\partial S_w}{\partial t} \tag{11}\]

Here \(S_w\) is the water saturation, \(\sigma_g'\) is the gravitational
effective overburden, \(P_c^{\mathrm{eq}}\) is the equilibrium capillary
pressure, and the final term is the rate-dependent dynamic capillary
stress of Hassanizadeh and Gray (1990) and Hassanizadeh et al.~(2002).
Sign conventions differ between soil-physics and rock-mechanics
formulations; the important point for \(\delta v/v\) is that the dynamic
coefficient \(\tau\) differs between wetting and drying, producing
hysteresis. During rapid drying, dynamic capillary suction can enhance
effective stress and stiffen the soil, while during wetting,
near-surface saturation and pore-pressure redistribution can soften it.

Shi et al.~(2026) showed three key results relevant to our framework:

\begin{enumerate}
\def\labelenumi{\arabic{enumi}.}
\item
  \textbf{Dynamic capillarity is essential.} Models with static
  capillary stress underestimate velocity increases during evaporation
  by \textasciitilde{}50\%, while models without any capillary stress
  produce significantly smaller velocity variations in both wetting and
  drying.
\item
  \textbf{Hysteresis is diagnostic.} The distinct wetting and drying
  \(\tau(S_w)\) coefficients produce \(\delta v/v\) hysteresis loops
  that are absent in standard poroelastic models. This hysteresis in the
  \(\delta v/v\)--saturation relationship is analogous to, but
  mechanistically distinct from, the mesoscopic nonlinearity (slow
  dynamics) observed in earthquake damage recovery (Snieder et al.,
  2017).
\item
  \textbf{Soil structure matters.} Tillage and compaction alter the pore
  network, modifying capillary behavior and drainage timescales.
  High-disturbance soils show pronounced evaporation-driven velocity
  rebounds and prolonged drainage, while undisturbed soils show minimal
  capillary effects due to efficient infiltration through well-connected
  pore networks.
\end{enumerate}

For our unified framework, the capillary effect extends Equation 4 to
partially saturated conditions. Above the water table, the effective
pressure becomes saturation-dependent, and the \(\delta v/v\) response
to rainfall involves not just pore-pressure diffusion but also the
dynamic redistribution of water between micropores and macropores. This
is particularly relevant for high-frequency \(\delta v/v\) measurements
(\textgreater{} 2 Hz) that sense the shallowest layers where saturation
changes are largest (Oakley et al., 2021). The workflow and regime map
(Figure 1; detailed frequency-depth domains in Figure S12) should
therefore include a ``capillary/vadose zone'' domain at high frequencies
and shallow depths, distinct from the saturated hydrological regime.

\begin{center}\rule{0.5\linewidth}{0.5pt}\end{center}

\section{5. Nonlinear Elasticity and the Acoustoelastic
Effect}\label{nonlinear-elasticity-and-the-acoustoelastic-effect}

\subsection{5.1 The Acoustoelastic Parameter and Its
Detection}\label{the-acoustoelastic-parameter-and-its-detection}

The Murnaghan (1937) equation of state under hydrostatic pressure ---
\(p = af + bf^2\) where \(a = 3\lambda + 2\mu\) and
\(b = 15\lambda + 10\mu - 27l - 9m - n\) (Figure S8) --- provides the
foundation for understanding velocity--pressure relations. The
acoustoelastic parameter \(\beta\) captures the sensitivity magnitude,
spanning three orders across materials (Figure 4; Clements and Denolle,
2023).

The \(\delta v/v\)--strain crossplot provides a diagnostic for
nonlinearity (Figure 6; detailed examples in Figure S7): linear confirms
third-order sufficiency, curvature indicates higher-order terms,
elliptical trajectories indicate viscoelastic phase lag, and hysteresis
indicates slow dynamics. Sens-Schönfelder and Eulenfeld (2019)
demonstrated that Earth tides probe in-situ nonlinearity at
\textasciitilde{}50 nanostrain amplitude.

\subsection{\texorpdfstring{5.2 Geological and Material Controls on
\(\beta\), \(\mu'\), and Nonlinear
Sensitivity}{5.2 Geological and Material Controls on \textbackslash beta, \textbackslash mu\textquotesingle, and Nonlinear Sensitivity}}\label{geological-and-material-controls-on-beta-mu-and-nonlinear-sensitivity}

The nonlinear elastic parameters are not fundamental material constants
--- they depend on rock \emph{microstructure}:

\textbf{Compliant porosity and microcracks.} The dominant nonlinearity
source in natural rocks is compliant porosity: thin, crack-like voids
with high aspect ratios that close under compression (Walsh, 1965;
Shapiro, 2003). The key parameters are crack density
\(\epsilon_c = N\langle a^3\rangle/V\) (Hudson, 1981), aspect ratio
\(\xi = w/a\) governing the closure pressure, and orientation
distribution controlling anisotropic nonlinearity (Sayers and Kachanov,
1995). Critically, \(|\beta|\) and \(\mu'\) \emph{decrease with
confining pressure} as cracks close, explaining why \(\delta v/v\)
sensitivity is strongest in the shallowest layers. Fokker et al.~(2021,
Fig. 2g) showed \(\mu' > 50\) for unconsolidated sediments but
\(\mu' < 10\) for consolidated rock.

\textbf{Grain contacts and cementation.} In unconsolidated granular
media, Hertz-Mindlin contact theory gives
\(\mu_{\text{eff}} \propto P_e^{1/3}\), producing
\(\mu' = d\mu/dP \sim \mu/P\) --- very large at low confining pressures
(Dvorkin and Nur, 1996). This explains why seasonal \(\delta v/v\) is
typically \textgreater{}0.1\% in sediments but \textless{}0.01\% in
crystalline rock.

\textbf{Mineralogy and cementation type.} Clay content increases
compliance and \(|\beta|\); salt cementation produces anomalously high
\(\partial(\rho v^2)/\partial\sigma_c \sim 1000\) at PATCX (Richter et
al., 2014) because salt bridges are highly stress-sensitive. Foliation
and bedding create intrinsic anisotropy that interacts with
stress-induced anisotropy.

\textbf{Fluid saturation state.} In fully saturated media, \(V_S\) is
insensitive to fluid type but sensitive to effective stress. In
partially saturated media, capillary suction stiffens the frame (Shi et
al., 2026); \(|\beta|\) can increase by factors of 2--5. Gas saturation
dramatically reduces \(V_P\) even at small gas fractions (patchy
saturation effects).

\begin{center}\rule{0.5\linewidth}{0.5pt}\end{center}

\section{6. Stress-Induced Anisotropy}\label{stress-induced-anisotropy}

\subsection{6.1 Directional Velocity Changes from Deviatoric
Stress}\label{directional-velocity-changes-from-deviatoric-stress}

The Tromp and Trampert (2018) framework (Eq. 4) shows that deviatoric
stress produces directionally dependent velocity changes. For a uniaxial
vertical stress \(T_{33}^0\), the induced deviatoric stress is
\(\boldsymbol{\tau}^0 = -(T_{33}^0/3)\,\text{diag}(1,1,-2)\).
Substituting into Equation 4 and evaluating for two end-member
propagation geometries --- vertically propagating S-waves
(\(\hat{\mathbf{k}} = \hat{\mathbf{z}}\),
\(\hat{\mathbf{a}} = \hat{\mathbf{x}}\)) versus horizontally propagating
S-waves with vertical polarization
(\(\hat{\mathbf{k}} = \hat{\mathbf{x}}\),
\(\hat{\mathbf{a}} = \hat{\mathbf{z}}\)) --- yields a velocity-change
difference between these two propagation directions:

\[\Delta\left(\frac{\delta V_S}{V_S}\right)_{\text{vert. vs. horiz.}} = \frac{T_{33}^0}{2\mu} \tag{*}\]

This directional splitting (Fokker et al., 2021, their Eqs. 9--11) is
proportional to the deviatoric stress divided by the shear modulus, and
is independent of \(\mu'\) --- meaning it exists even without nonlinear
elasticity, arising purely from the stress-modified constitutive
relation (Figure 5). For a fixed propagation direction, the SV-SH
splitting (velocity difference between vertically and horizontally
polarized shear waves) is likewise proportional to the deviatoric stress
but depends on the specific geometry of propagation relative to the
stress axis.

\subsection{6.2 Microcrack Closure and the Parkfield
Observation}\label{microcrack-closure-and-the-parkfield-observation}

At Parkfield, Okubo et al.~(2024) found that the long-term
\(\delta v/v\) increase correlates with the \emph{contractional} axial
strain (oriented N35°W to N45°E) but \emph{not} with the dilatational
strain, which shows a slight extension. This anisotropic response is
naturally explained by stress-induced microcrack closure (Sayers and
Kachanov, 1995; Verdon et al., 2008): cracks oriented perpendicular to
the maximum compressive stress close preferentially, increasing rigidity
and velocity in that direction (Figure 5), while cracks parallel to the
compression remain open.

The scalar (isotropic) \(\delta v/v\) --- as typically measured from
coda waves averaging over all propagation directions --- integrates over
the azimuthal dependence. For a traceless deviatoric sensitivity kernel,
this average is zero, meaning the isotropic \(\delta v/v\) is blind to
deviatoric stress (Tromp and Trampert, 2018). Detecting the deviatoric
component therefore requires either \emph{azimuthal binning} of
\(\delta v/v\) (measuring velocity change as a function of inter-station
azimuth) or \emph{multi-component analysis} (separating Rayleigh and
Love wave contributions, as in Fokker et al., 2021). This is an
underexploited observational strategy that could provide constraints on
stress orientation from passive seismic monitoring.

\begin{center}\rule{0.5\linewidth}{0.5pt}\end{center}

\section{7. Rheological Models and Depth-Resolved Joint
Inversion}\label{rheological-models-and-depth-resolved-joint-inversion}

\subsection{7.1 Post-Seismic Healing}\label{post-seismic-healing}

The logarithmic healing model (Snieder et al., 2017; Okubo et al.,
2024):

\[L(t) = -\int_{\tau_{\min}}^{\tau_{\max}} \frac{1}{\tau} \exp\left(-\frac{t - t_{EQ}}{\tau}\right) d\tau \tag{12}\]

encodes the distribution of relaxation timescales. At Parkfield,
\(\tau_{\max}\) is constrained to 1--30,000 years, meaning post-2004
healing may continue today (Figure 6; detailed healing curves in Figure
S9). The trade-off between healing and tectonic trend is a fundamental
limitation requiring independent geodetic constraints.

\subsection{7.2 Distinguishing Rheological
Models}\label{distinguishing-rheological-models}

Different rheologies produce distinct \(\delta v/v\)--strain signatures
(Figure 6): elastic (linear crossplot), Maxwell (exponential decay),
Kelvin-Voigt (delayed buildup), SLS (frequency-dependent dispersion,
elliptical crossplot), and slow dynamics (logarithmic recovery,
hysteresis). Tidal modulation (Sens-Schönfelder and Eulenfeld, 2019)
provides the cleanest diagnostic at well-known forcing periods.

\subsection{7.3 Frequency-Dependent Depth Tomography of
Stress}\label{frequency-dependent-depth-tomography-of-stress}

A key opportunity that has been underexploited is the \emph{frequency
dependence} of \(\delta v/v\). Surface waves at different frequencies
have sensitivity kernels peaking at different depths (Obermann et al.,
2014):

\[\frac{\delta v}{v}(f, t) = \int_0^\infty K(z, f) \cdot \frac{\delta V_S}{V_S}(z, t) \, dz \tag{13}\]

where \(K(z, f)\) is the Rayleigh or Love wave phase-velocity
sensitivity kernel at frequency \(f\) (Figure 2a). The peak sensitivity
depth is approximately \(V_s / (3f)\), ranging from \textasciitilde{}50
m at 2 Hz to \textasciitilde{}5 km at 0.1 Hz for typical crustal
velocities.

By measuring \(\delta v/v\) simultaneously at multiple frequency bands
(e.g., 0.1--0.5, 0.5--1.0, 1.0--2.0, 2.0--4.0 Hz, as in Okubo et al.,
2024 and Ermert et al., 2023), one obtains a set of equations --- one
per frequency band --- each sampling a different depth range. This is
formally a linear inverse problem:

\[\mathbf{d} = \mathbf{K} \mathbf{m} \tag{14}\]

where \(\mathbf{d}\) is the vector of \(\delta v/v\) at \(N_f\)
frequencies, \(\mathbf{K}\) is the kernel matrix (dimensions
\(N_f \times N_z\)), and \(\mathbf{m}\) is the depth profile of
\(\delta V_S / V_S\), which through the nonlinear elastic relation (Eq.
3) maps to the depth profile of stress or strain perturbation.

This inverse problem is strongly ill-conditioned: surface-wave and coda
sensitivity kernels are broad, adjacent frequency bands are correlated,
and the material sensitivity \(\beta(z)\) or \(\mu'(z)\) is rarely known
independently. Any depth-resolved result therefore requires
regularization, uncertainty propagation, and resolution tests. In this
paper, depth-dependent inversion is a proposed workflow rather than a
demonstrated result.

\subsection{7.4 Lapse-Time Windows as Part of the Sensitivity
Kernel}\label{lapse-time-windows-as-part-of-the-sensitivity-kernel}

The observable is not simply \(\delta v/v(f,t)\). It is better written
as

\[d(f,\tau,W,T_{\mathrm{stack}},t) = \int_0^\infty K(z; f,\tau,W,\mathbf{g}) \, \frac{\delta V_S}{V_S}(z,t) \, dz + \epsilon \tag{15}\]

where \(f\) is the frequency band, \(\tau\) is the coda-window center
lapse time, \(W\) is the window duration, \(T_{\mathrm{stack}}\) is the
substack length, \(\mathbf{g}\) describes source-receiver geometry and
wavefield directionality, and \(\epsilon\) includes measurement error
and source-side bias. This notation makes explicit that coda-window
choice changes the kernel, not merely the uncertainty.

Takano et al.~(2019) demonstrated this directly at Izu-Oshima volcano by
estimating tidal-strain sensitivity from ambient-noise
cross-correlations at different lapse times. In their 2--4 Hz analysis,
early lapse windows (2--7 s) showed strong velocity sensitivity to tidal
dilatation and contraction, while later windows (7--35 s) showed reduced
strain sensitivity. Array analysis indicated apparent velocities close
to the local Rayleigh-wave group velocity in the early windows and
higher apparent velocities in later windows, consistent with scattered
or reflected body-wave contributions from greater depth. The result is
not that early windows are universally best, but that lapse time is a
physical axis of the measurement.

This creates a reproducibility problem and an opportunity. If different
studies choose windows by convention, visual inspection, or inherited
defaults, their \(\delta v/v\) estimates may differ because they sample
different wave types, depths, and nonlinear sensitivities. A
reproducible workflow should therefore compute a rolling lapse-time
profile from early to late coda and score each overlapping window using
explicit metrics:

\[J = \sum_i w_i Q_i(f,\tau,W,T_{\mathrm{stack}}) \tag{16}\]

where \(Q_i\) are normalized scores and \(w_i\) are analysis-specific
weights. The core metrics should include:

\begin{enumerate}
\def\labelenumi{\arabic{enumi}.}
\tightlist
\item
  \textbf{Measurement quality:} waveform coherence, signal-to-noise
  ratio, stretching or MWCS fit sharpness, and estimated \(\delta v/v\)
  uncertainty.
\item
  \textbf{Wave-type consistency:} apparent velocity or array slowness
  compared with expected Rayleigh/Love group velocities, flagging
  windows likely dominated by body waves or mixed wave types.
\item
  \textbf{Depth targeting:} overlap between \(K(z; f,\tau,W)\) and the
  hypothesized target depth interval, or a frequency-depth proxy when
  full kernels are not yet available.
\item
  \textbf{Lapse-time stability:} sensitivity of the recovered
  \(\delta v/v\) or \(\beta\) estimate to adjacent windows, used to
  identify transitions in wave type or scattering regime.
\item
  \textbf{Temporal resolution:} substack length short enough to resolve
  the forcing period but long enough to maintain measurement precision.
\item
  \textbf{Source stability:} robustness to changes in noise-source
  spectrum, azimuth, and seasonal source distribution.
\end{enumerate}

The future \texttt{codameter} package should implement this as a rolling
window-selection layer before mechanism inversion. Rather than reporting
a single hand-picked coda window, a study should report the
early-to-late lapse profile, the metric weights, the stable intervals or
transitions, and the sensitivity of the scientific conclusion to that
ensemble. This would turn an under-documented processing choice into a
reproducible part of the physical forward model.

A Bayesian framing of this workflow treats each window
\(M_j=(f_j,\tau_j,W_j,T_{\mathrm{stack},j})\) as a candidate measurement
and asks how much a scientific target \(\theta\) --- a \(\delta v/v\)
value, tidal strain sensitivity, or hydrological regression coefficient
--- depends on that choice. If each window yields an estimate with
posterior \(p(\theta|y,M_j)\) and a weight \(p(M_j|y)\), a
model-averaged posterior and its variance can be written

\[p(\theta|y)=\sum_j p(M_j|y)\,p(\theta|y,M_j), \tag{17}\]

\[\mathrm{Var}(\theta|y)=\underbrace{\mathbb{E}_{M}[\mathrm{Var}(\theta|y,M)]}_{\text{within-window (measurement)}}+\underbrace{\mathrm{Var}_{M}[\mathbb{E}(\theta|y,M)]}_{\text{between-window (method)}}. \tag{18}\]

Equation 18 is the familiar within-group/between-group variance
decomposition: the first term is the average measurement uncertainty
inside a fixed window, and the second is the spread of the per-window
point estimates. We stress that Equation 18 should be used as a
\emph{window-sensitivity diagnostic} --- a check on whether a conclusion
survives plausible processing choices --- and not as a turnkey inference
method, for two reasons. First, heuristic weights of the form
\(p(M_j|y)\propto\exp(\lambda J_j)\) built from the quality score
\(J_j\) are not Bayesian model evidence; the resulting variance depends
on tuning constants (the temperature \(\lambda\) and the metric
normalizations) rather than on the data, so defensible weights require a
genuine marginal likelihood or out-of-sample predictive score. Second,
and more fundamentally, overlapping lapse windows share coda samples
and, by Equation 15, measure \emph{physically different} depth-weighted
averages; a large between-window term therefore often reflects resolved
depth and wave-type structure rather than method error. The principled
route is consequently not to average windows but to invert them jointly:
treat the windowed measurements \(\{d_j\}\) as correlated data with
known kernels \(K(z;f,\tau,W)\) (Eq. 15) and let the posterior on the
depth field \(\delta V_S/V_S(z,t)\) carry the resolution uncertainty
(Section 7.5). In that hierarchical formulation the window metrics enter
as a heteroscedastic noise model and quality gate rather than as model
weights, and the decomposition in Equation 18 reappears in its correct
form as the data-noise versus posterior-resolution split. Here we use
this only as a window-sensitivity \emph{diagnostic}; the full
reproducibility and uncertainty-quantification treatment --- a
hierarchical estimator with a correlated-window data covariance,
quality-as-noise weighting, lapse-time regime segmentation, and
coverage-validated credible intervals --- is developed in a companion
methods paper and the \texttt{codameter} software, because it requires
its own synthetic-recovery and multi-station real-data validation
distinct from the interpretation focus of this work.

\subsection{7.5 Toward 3-D Stress/Strain
Imaging}\label{toward-3-d-stressstrain-imaging}

When combined with geodetic observations, the system becomes better
constrained, although still non-unique:

\textbf{Vertical resolution} comes from the frequency dependence of
\(\delta v/v\) (different frequencies → different depths).

\textbf{Horizontal resolution} comes from the spatial array geometry:
\(\delta v/v\) measured at different station pairs samples different
lateral regions, while GNSS provides point measurements of surface
displacement and InSAR provides spatially continuous line-of-sight
deformation.

\textbf{Directional resolution} comes from multi-component
\(\delta v/v\) analysis (ZZ, RR, TT, RT cross-correlations), which
samples different combinations of Rayleigh and Love waves with different
sensitivities to isotropic and anisotropic velocity changes.

\textbf{Temporal resolution} is inherently different: geodetic
observations respond elastically (instantaneously to loading), while
\(\delta v/v\) integrates both elastic and inelastic effects. The
\emph{temporal mismatch} between geodetic strain and \(\delta v/v\) is
itself a diagnostic of rheology:

\begin{itemize}
\tightlist
\item
  If \(\delta v/v\) tracks geodetic strain instantaneously → elastic
  response
\item
  If \(\delta v/v\) lags behind geodetic strain → viscoelastic
  retardation
\item
  If \(\delta v/v\) shows logarithmic recovery while geodetic strain is
  steady → slow dynamics
\item
  If \(\delta v/v\) shows hysteresis with geodetic strain → dynamic
  capillary or mesoscopic nonlinear effects
\end{itemize}

The full joint inversion then seeks to determine the 4-D (x, y, z, t)
stress/strain field \(\sigma_{ij}(\mathbf{x}, t)\) that simultaneously
satisfies:

\begin{enumerate}
\def\labelenumi{\arabic{enumi}.}
\tightlist
\item
  \textbf{\(\delta v/v\) forward model} (Eq. 13): the
  frequency-dependent, depth-averaged velocity change at each station
\item
  \textbf{Geodetic forward model}: the surface displacement field from
  elastic/viscoelastic deformation
\item
  \textbf{Constitutive relation}: the mapping between stress/strain and
  velocity change via nonlinear elasticity (Eqs. 3--5, 7)
\item
  \textbf{Equilibrium constraint}: the induced stress must satisfy
  \(\nabla \cdot \boldsymbol{\tau}^0 = 0\) (Tromp and Trampert, 2018)
\item
  \textbf{Physical forcing models}: thermoelastic (Eqs. 8--9),
  hydrological (Eq. 5), tectonic (from GNSS velocities)
\end{enumerate}

The model parameters include depth-dependent material properties
(\(\beta(z)\), \(\mu'(z)\), \(\kappa_T\), \(c\)) and rheological
parameters (\(\tau_{\min}\), \(\tau_{\max}\), viscoelastic relaxation
time). The frequency dependence of \(\delta v/v\) is the key ingredient
that can convert a surface measurement into a depth-sensitive probe.
With sufficient bandwidth, array aperture, and external constraints,
this may support 3-D stress/strain models analogous to time-lapse
seismic tomography but operating continuously and passively.

This approach has concrete applications: in volcanic settings,
multi-frequency \(\delta v/v\) combined with GNSS tilt can distinguish
between shallow magma chamber pressurization (high-frequency
\(\delta v/v\) signal) and deep magma supply (low-frequency signal +
GNSS displacement). In hydrological settings, the depth-dependent
balance between loading and pore-pressure effects (Section 4.3) can be
resolved rather than assumed. In tectonic settings, the depth of stress
accumulation can be constrained, potentially distinguishing between
shallow fault locking and deeper aseismic creep.

\begin{center}\rule{0.5\linewidth}{0.5pt}\end{center}

\section{8. Sensitivity Analysis and
Validity}\label{sensitivity-analysis-and-validity}

\subsection{8.1 Homogeneous Half-Space}\label{homogeneous-half-space}

The homogeneous half-space assumption fails when velocity contrasts
exceed \textasciitilde{}2:1 within the sensitivity kernel (Figure 2;
detailed validity tests in Figure S10). At soft sedimentary sites ---
such as Mexico City, where \(V_S\) ranges from 50 m/s in the lake zone
to 800 m/s in bedrock --- the homogeneous assumption breaks down at all
frequencies, requiring station-specific velocity profiles and proper
sensitivity kernels (Ermert et al., 2023). At hard-rock stations above 1
Hz, sensitivity is confined to the upper few meters where the medium is
more likely to be approximately uniform.

\subsection{8.2 Linear
Acoustoelasticity}\label{linear-acoustoelasticity}

The linear acoustoelastic approximation is valid for strains below
\textasciitilde{}10\(^{-5}\) (Figure S11), covering tidal,
thermoelastic, and most hydrological signals, but breaks down for
coseismic and strong-motion strains. The regime diagram (Figure 1;
detailed version in Figure S12) maps dominant processes versus frequency
and depth, including the capillary/vadose zone regime at high
frequencies.

Table 1 summarizes the main parameters, ranges, and validity limits
needed for applying the framework; the full parameter overview is
provided as Table S1.

{\def\LTcaptype{none} % do not increment counter
\begin{longtable}[]{@{}
  >{\raggedright\arraybackslash}p{(\linewidth - 6\tabcolsep) * \real{0.2308}}
  >{\raggedleft\arraybackslash}p{(\linewidth - 6\tabcolsep) * \real{0.3077}}
  >{\raggedright\arraybackslash}p{(\linewidth - 6\tabcolsep) * \real{0.2308}}
  >{\raggedright\arraybackslash}p{(\linewidth - 6\tabcolsep) * \real{0.2308}}@{}}
\toprule\noalign{}
\begin{minipage}[b]{\linewidth}\raggedright
Parameter
\end{minipage} & \begin{minipage}[b]{\linewidth}\raggedleft
Typical range
\end{minipage} & \begin{minipage}[b]{\linewidth}\raggedright
Primary control
\end{minipage} & \begin{minipage}[b]{\linewidth}\raggedright
Main limitation
\end{minipage} \\
\midrule\noalign{}
\endhead
\bottomrule\noalign{}
\endlastfoot
\(\kappa_T\) & 0.15--2.0 mm\(^2\)/s & Thermoelastic skin depth &
Strongly site and moisture dependent \\
\(c\) & \(10^{-3}\)--\(10^2\) m\(^2\)/s & Pore-pressure diffusion &
Varies by orders of magnitude across lithology \\
\(\mathrm{Pe}=\omega L^2/c\) & \(\ll1\) to \(\gg1\) & Selects drained
(\(\kappa\)) vs undrained (\(\kappa_u\)) modulus in Eq. 7 & Data-driven
per signal; needs \(c\) and \(L\) \\
\(B\), \(\nu_u\) & 0.2--0.95, 0.25--0.49 & Drained versus undrained
response & Requires saturation and poroelastic constraints \\
\(\beta\) & \(10\)--\(10^4\) & Velocity-strain sensitivity & Directional
under deviatoric loading \\
\(\mu'\) & \(10\)--\(10^3\) & Stress-to-velocity conversion & Depends on
pressure, cracks, and sediment fabric \\
\(\tau_{\min}\), \(\tau_{\max}\) & days to \(>10^4\) yr & Slow-dynamics
recovery & Trades off with long-term tectonic trend \\
\(K(z,f)\) & site specific & Frequency-depth sensitivity & Broad,
correlated kernels make inversion ill-conditioned \\
\end{longtable}
}

\subsection{8.3 Spatial Generalization: From 1-D to
3-D}\label{spatial-generalization-from-1-d-to-3-d}

The forward models in Sections 3--6 assume laterally homogeneous (1-D)
structure. Generalizing to 3-D requires accounting for spatially varying
velocity structure, material properties, and observation geometry.

\textbf{Spatially varying sensitivity kernels.} The kernel \(K(z, f)\)
in Equation 1 depends on the local \(V_S(z)\), \(V_P(z)\), and
\(\rho(z)\) profiles. At sites with laterally varying geology ---
sedimentary basins (Ermert et al., 2023), fault zones (Hillers et al.,
2015), volcanic edifices (Donaldson et al., 2019) --- the kernel differs
between station pairs. Ermert et al.~(2023) addressed this by computing
station-specific 1-D profiles using geotechnical classifications and
aquitard thickness. A fully 3-D treatment requires 3-D sensitivity
kernels computed with adjoint methods (Tromp et al., 2005) or
spectral-element simulations.

\textbf{Spatially varying material properties.} The nonlinear parameters
(\(\beta\), \(\mu'\), \(\kappa_T\), \(c\)) vary laterally. In a
sedimentary basin, \(\mu'\) is large in the basin fill and small in the
basement; hydraulic diffusivity \(c\) can vary by orders of magnitude
across fault zones and clay layers. Clements and Denolle (2023)
demonstrated this: basin sites show low diffusivity (long hydrological
memory) while mountain sites show high diffusivity (rapid drainage).

\textbf{What additional Earth observations are needed?} Generalizing to
3-D requires: (1) a \textbf{3-D \(V_S\) model} from ambient noise
tomography at the monitoring frequencies, providing the base model for
sensitivity kernels; (2) a \textbf{\(V_P/V_S\) ratio} from joint
Rayleigh-Love inversion, constraining Poisson's ratio for the
thermoelastic model; (3) a \textbf{density model} from gravity or
empirical relations, for converting \(\delta V_S/V_S\) to
\(\delta\mu/\mu\); (4) \textbf{near-surface geotechnical data}
(\(V_{S30}\), water table depth, porosity) constraining \(\mu'\),
\(\kappa_T\), and saturation state at each site; (5) \textbf{co-located
geodetic coverage} (GNSS, InSAR) for independent strain constraints; and
(6) \textbf{distributed meteorological and hydrological data} co-located
with the seismic array. Dense arrays --- nodal deployments, distributed
acoustic sensing (DAS; Rodríguez Tribaldos and Ajo-Franklin, 2021; Shi
et al., 2026) --- provide the spatial resolution needed to resolve
lateral variations in both \(\delta v/v\) and the material properties
controlling it.

\begin{center}\rule{0.5\linewidth}{0.5pt}\end{center}

\section{9. Application to Field
Observations}\label{application-to-field-observations}

\emph{{[}Data application --- placeholder.{]}} A quantitative
application of the framework to field observations is in preparation. A
preliminary reinterpretation of three contrasting published datasets ---
Parkfield (strike-slip; Okubo et al., 2024), Cascadia (subduction;
Kidiwela et al., 2026), and Kīlauea (volcanic caldera collapse;
Hotovec-Ellis et al., 2022) --- is provided in Text S1 of the Supporting
Information (with its three-site synthesis figure and comparison table).
That preliminary analysis uses only published \(\delta v/v\)
measurements, velocity profiles, and geodetic strain constraints (no new
waveform processing) to illustrate how the bridge relation (Eq. 7), the
data-driven drainage regime, and the isotropic-versus-deviatoric
diagnostic operate across settings. A full application with
independently reprocessed \(\delta v/v\) records is deferred to future
work (§10.5).

\section{10. Discussion}\label{discussion}

\subsection{10.1 The Unifying Power of Nonlinear
Elasticity}\label{the-unifying-power-of-nonlinear-elasticity}

All \(\delta v/v\) sources --- thermoelastic, hydrological, tectonic,
volcanic --- operate through the same mechanism: stress-dependent
modification of elastic moduli via compliant porosity and microcracks.
The \emph{same} material parameters (\(\beta\), \(\mu'\),
\(\partial(\rho v^2)/\partial\sigma_c\)) should consistently explain all
forcings at a given site. Inconsistencies diagnose either unmodeled
forcings or violations of the isotropic assumption.

\subsection{10.2 The Capillary Frontier}\label{the-capillary-frontier}

The Shi et al.~(2026) results on dynamic capillary effects represent a
significant extension of the \(\delta v/v\) framework into unsaturated
media. The hysteretic \(\delta v/v\) response to wetting-drying cycles
--- governed by the rate-dependent dynamic coefficient \(\tau(S_w)\) ---
is mechanistically distinct from both the Fokker et al.~(2021)
poroelastic framework (which assumes full saturation) and the Snieder et
al.~(2017) slow dynamics (which involves crack healing). This opens a
``capillary window'' for monitoring soil moisture, agricultural
disturbance, and vadose-zone hydrology with seismic methods. Future work
should integrate the Hertz-Mindlin + dynamic capillarity model of Shi et
al.~(2026) into the depth-integrated forward model (Eq. 1) to predict
multi-frequency \(\delta v/v\) in partially saturated settings.

\subsection{10.3 Implications for Monitoring
Applications}\label{implications-for-monitoring-applications}

For \textbf{volcanic monitoring}, environmental signals of
\textasciitilde{}0.05\% can mask pre-eruptive signals. The framework
enables forward modeling of expected environmental \(\delta v/v\), and
the anisotropic formulation can detect oriented magmatic stress. The
Kīlauea case (§9.3) demonstrates that the dominant fracture fabric ---
ring fractures vs.~radial dikes --- determines the \emph{sign} of the
\(\delta v/v\) response to reservoir pressurization, a prediction
testable at other caldera systems.

For \textbf{tectonic monitoring}, the Parkfield case shows
\(\delta v/v\) detects inter-seismic loading at
\textasciitilde{}0.005\%/yr. The deviatoric component is essential ---
dilatational strain alone cannot explain the observation (Okubo et al.,
2024).

For \textbf{hydrological monitoring}, the depth-resolved approach
(Section 7.3) can separate shallow capillary effects from deep
pore-pressure changes, and the competition between loading and
pore-pressure effects can be resolved rather than assumed.

\subsection{10.4 Practical Workflow for Applying the Framework at a New
Site}\label{practical-workflow-for-applying-the-framework-at-a-new-site}

We summarize the steps a researcher would take to apply this framework:

\begin{enumerate}
\def\labelenumi{\arabic{enumi}.}
\tightlist
\item
  \textbf{Characterize the site.} Obtain a 1-D (or 3-D) \(V_S(z)\),
  \(V_P(z)\), and \(\rho(z)\) profile from ambient noise tomography,
  active surveys, or geotechnical databases. Estimate \(\mu(z)\) and
  \(\kappa(z)\) from these profiles.
\item
  \textbf{Design and score measurement windows.} Define a candidate grid
  of frequency bands, coda-window starts and ends, and substack lengths.
  Rank windows using coherence, signal-to-noise ratio, uncertainty,
  apparent wave type, depth targeting, lapse-time stability, and
  source-stability metrics (§7.4). Report the candidate grid and the
  selected-window ensemble.
\item
  \textbf{Compute sensitivity kernels.} Use the velocity profile to
  compute Rayleigh and Love wave phase-velocity sensitivity kernels
  \(K(z,f)\) at the monitoring frequencies (e.g., using the adjoint
  method of Tromp et al., 2005, or eigenfunction codes). Where possible,
  extend the kernel description to include lapse time and window
  duration, \(K(z;f,\tau,W)\).
\item
  \textbf{Estimate nonlinear parameters.} Compute \(\mu'(z) = d\mu/dP\)
  from the shear-modulus--pressure profile, or use empirical values for
  the local lithology. The bridge relation (Eq. 7) then gives
  \(\beta(z)\).
\item
  \textbf{Forward model environmental \(\delta v/v\).} Using local
  temperature records, compute the thermoelastic \(\delta v/v\) via the
  Berger-Richter model (Eqs. 8--9). Using precipitation and/or
  groundwater data, compute the hydrological \(\delta v/v\) via the
  Roeloffs-Fokker model (Eq. 5). Integrate depth-dependent
  \(\delta V_S/V_S(z,t)\) through the kernels (Eq. 1 or Eq. 15) to
  predict \(\delta v/v(f,\tau,W,t)\).
\item
  \textbf{Subtract and interpret residuals.} Remove the modeled
  environmental contributions from the observed \(\delta v/v\) time
  series. Residuals may contain tectonic, volcanic, or anthropogenic
  signals. A robust residual should persist across the top-ranked window
  ensemble.
\item
  \textbf{Compare with geodesy.} Overlay residual \(\delta v/v\) with
  GNSS-derived strain or InSAR deformation. The temporal relationship
  (instantaneous tracking, lag, or hysteresis) diagnoses the rheological
  regime (§7.2). Multi-frequency and multi-lapse residuals constrain the
  depth distribution and wave-type dependence of the anomalous signal
  (§7.3--7.4).
\end{enumerate}

\subsection{10.5 Limitations}\label{limitations}

The data applications in Section 9 use published \(\delta v/v\)
measurements and velocity profiles rather than reprocessed waveform
data; the quantitative predictions (stress rates, \(\beta\), \(\mu'\))
depend on the accuracy of the input values from Okubo et al.~(2024),
Kidiwela et al.~(2026), Hotovec-Ellis et al.~(2022), and the respective
velocity models. The Parkfield GNSS cross-check has a factor-of-1.4
discrepancy that may reflect depth-transfer effects between the surface
strain measurement and the 0.8 km sensitivity depth. The Cascadia
borehole comparison is a calibration-based consistency check rather than
an independent validation because the borehole strain is used to
determine \(\beta\). The Kīlauea estimate uses a simplified spheroidal
reservoir geometry; more complex geometries (including the piston and
ring fault) would modify the strain field. The directional bridge
relation (\(\mu'\) from \(\beta_{\text{axial}}\) or
\(\beta_{\text{radial}}\)) is an order-of-magnitude approximation;
Equation 7 was derived for isotropic loading.

The companion notebook figures use \emph{synthetic} \(\delta v/v\) with
physically realistic shapes to illustrate the framework's forward
predictions and diagnostic signatures; they are not fits to field data.
Quantitative data applications in this paper rest on the three published
studies (Okubo et al., 2024; Kidiwela et al., 2026; Hotovec-Ellis et
al., 2022). End-to-end validation against independently reprocessed
\(\delta v/v\) records --- selecting several stations with public
waveform data, measuring \(\delta v/v\), and applying the full
forward-and-subtract workflow of §10.4 --- is a clear next step and is
left to future work.

Additional limitations include: the forward models are primarily 1-D,
while real geology has lateral heterogeneity (§8.3). Source-side
noise-field variations can produce apparent \(\delta v/v\) (Zhan et al.,
2013; Okubo et al., 2024 mitigate this with multi-component analysis).
The decomposition into forcings (Eq. 6) is non-unique without
independent geodetic constraints. Mesoscopic nonlinearity at high
strains requires constitutive models beyond the Murnaghan framework
(Gassenmeier et al., 2016). The capillary model (§4.4) is presented as
external evidence from Shi et al.~(2026) rather than implemented in the
companion notebooks. The proposed lapse-window objective (§7.4) is a
design framework, not yet a universal standard; its weights should be
tuned to the scientific target and validated against synthetic recovery
tests and independent strain or hydrological observations.

\begin{center}\rule{0.5\linewidth}{0.5pt}\end{center}

\section{11. Conclusions}\label{conclusions}

\begin{enumerate}
\def\labelenumi{\arabic{enumi}.}
\tightlist
\item
  All major sources of \(\delta v/v\) variability operate through
  nonlinear elasticity. The Murnaghan (1937) framework provides the
  formal foundation.
\item
  The strain formulation (\(\delta v/v = \beta \epsilon_{kk}\)) and
  stress formulation (\(\delta V_S/V_S = \mu' p^0 / 2\mu\)) are unified
  under isotropic loading by the bridge relation
  \(\beta = -\mu'\kappa/(2\mu)\) (Eq. 7), but diverge under deviatoric
  loading --- a divergence that carries diagnostic information about
  stress orientation.
\item
  Thermoelastic \(\delta v/v\) is controlled primarily by
  \(\partial(\rho v^2)/\partial\sigma_c\) (50--1000), with annual
  amplitudes of 0.01--0.3\%.
\item
  Hydrological loading produces competing effects whose balance depends
  on permeability and drainage (Fokker et al., 2021).
\item
  Dynamic capillary effects in partially saturated media produce
  hysteretic \(\delta v/v\) distinct from standard poroelastic responses
  (Shi et al., 2026).
\item
  Stress-induced anisotropy from microcrack closure explains why
  \(\delta v/v\) correlates with contractional but not dilatational
  strain (Okubo et al., 2024). Detecting deviatoric stress requires
  azimuthal binning or multi-component analysis.
\item
  At Parkfield, the framework predicts a deviatoric stress rate of
  \textasciitilde12 kPa/yr at \textasciitilde0.8 km depth with
  \(\beta_{\text{axial}} \approx 240\) and \(\mu' \approx 250\),
  cross-checked against GNSS strain. At Cascadia, the
  borehole-calibrated framework gives an isotropic stress rate of 0.58
  kPa/yr at \textasciitilde0.2 km depth with \(\beta \approx 3160\) and
  \(\mu' \approx 620\). At Kīlauea, co-collapse radial stress of
  \textasciitilde170 kPa is recovered from \(\delta v/v \approx 0.5\%\)
  with \(\beta_{\text{radial}} \approx 300\) and \(\mu' \approx 360\),
  consistent with the geodetically constrained reservoir pressure change
  within a factor of 2. The effective \(|\beta|\) for fractured
  crystalline rock (granite, basalt) clusters around 200--400, while
  unconsolidated marine sediment gives \(|\beta| \sim 3000\).
\item
  The loading geometry and fracture fabric jointly determine which
  framework applies: isotropic \(\delta v/v = \beta\epsilon_{kk}\)
  succeeds at Cascadia (volumetric compression, no preferred fracture
  orientation) but fails at both Parkfield (strike-slip shear,
  fault-parallel cracks) and Kīlauea (magmatic pressurization,
  concentric ring fractures). The dominant fracture orientation
  determines the controlling strain component --- fault-normal
  contraction at Parkfield, radial compression at Kīlauea --- and this
  directional sensitivity explains the opposite \(\delta v/v\) response
  to pressurization observed at Piton de la Fournaise (radial dike
  fractures).
\item
  Multi-frequency \(\delta v/v\) combined with GNSS/InSAR has the
  potential to enable depth-resolved 3-D stress/strain imaging, provided
  sensitivity kernels, regularization, and independent material
  constraints are available. At Cascadia, the 2:1 ratio of low-frequency
  \(\delta v/v\) trends between northern and central segments is
  consistent with the locking contrast inferred from onshore geodesy
  under comparable-kernel and comparable-\(\beta\) assumptions.
\item
  Coda-window choice is a physical part of the measurement, not a
  neutral processing detail. The observable should be reported as
  \(\delta v/v(f,\tau,W,T_{\mathrm{stack}},t)\), and window selection
  should be based on explicit metrics for coherence, uncertainty, wave
  type, depth targeting, lapse-time stability, temporal resolution, and
  source stability. This is the first target for the \texttt{codameter}
  implementation because it directly addresses reproducibility; the full
  hierarchical uncertainty quantification of dv/v with respect to these
  processing choices is developed in a companion methods paper.
\item
  Rheological diagnostics from \(\delta v/v\)--strain crossplots and
  tidal modulation distinguish elastic, viscoelastic, and slow-dynamics
  behavior. At Parkfield, a dual-population model (elastic tectonic +
  slow-dynamics healing) is required.
\item
  Linear acoustoelasticity is valid for strains below
  \textasciitilde{}10\(^{-5}\); the nonlinear parameters (\(\beta\),
  \(\mu'\)) are controlled by rock microstructure and decrease with
  confining pressure, spanning \(|\beta| \sim 240\) (fractured granite)
  to \(|\beta| \sim 3160\) (marine sediment).
\end{enumerate}

\begin{center}\rule{0.5\linewidth}{0.5pt}\end{center}

\section{Statement of AI Use}\label{statement-of-ai-use}

The sole author of this work is M. A. Denolle, who is responsible for
all scientific content, claims, and conclusions. AI tools (Claude,
Anthropic) were used as a research and writing assistant under the
author's direction --- for literature synthesis, numerical
implementation, figure generation, and drafting --- analogous to the use
of conventional computational and editorial tools. The author conceived
the study, provided the scientific direction and problem framing,
selected the knowledge base, contributed the key insight to integrate
partially saturated media with dynamic capillary effects, and verified,
edited, and takes responsibility for the final text and all results. All
scientific claims are grounded in published, peer-reviewed literature,
with numerical inputs traced to their sources in
\texttt{docs/site\_analyses/provenance\_tables.md}. Full documentation
of the prompts, AI reasoning, and model information is provided in
\texttt{docs/ai\_documentation/}, consistent with AGU's policy that AI
tools are not eligible for authorship.

\begin{center}\rule{0.5\linewidth}{0.5pt}\end{center}

\section{Acknowledgments}\label{acknowledgments}

M.A.D. acknowledges support from the Packard Foundation and the Murdock
Charitable Trust. The project knowledge base includes papers by Okubo et
al.~(2024), Clements and Denolle (2023), Ermert et al.~(2023), Richter
et al.~(2014), Fokker et al.~(2021), Tromp and Trampert (2018),
Murnaghan (1937), Shi et al.~(2026), and Verdon (2008). Additional
literature and metadata checks were performed with disclosed AI-assisted
search tools and verified against cited sources where used.

\section{Data Availability}\label{data-availability}

The companion Jupyter notebooks and figure-generation code are
maintained in this repository and will be archived on Zenodo before
publication; the archive DOI is pending. The notebooks use standard
Python scientific libraries (NumPy, SciPy, Matplotlib).

\begin{center}\rule{0.5\linewidth}{0.5pt}\end{center}

\begin{center}\rule{0.5\linewidth}{0.5pt}\end{center}

\section{References}\label{references}

Ben-Zion, Y., \& Leary, P. (1986). Thermoelastic strain in a half-space
covered by unconsolidated material. \emph{Bulletin of the Seismological
Society of America}, 76, 1447--1460.

Anderson, K. R., Johanson, I. A., Patrick, M. R., Gu, M., Segall, P.,
Poland, M. P., et al.~(2019). Magma reservoir failure and the onset of
caldera collapse at Kīlauea volcano in 2018. \emph{Science}, 366,
eaaz1822.

Boness, N. L., \& Zoback, M. D. (2006). Mapping stress and structurally
controlled crustal shear velocity anisotropy with a multi-layered model.
\emph{Geophysical Research Letters}, 33, L01304.

Berger, J. (1975). A note on thermoelastic strains and tilts.
\emph{Journal of Geophysical Research}, 80, 274--277.

Birch, F. (1961). The velocity of compressional waves in rocks to 10
kilobars, part 2. \emph{Journal of Geophysical Research}, 66,
2199--2224.

Brenguier, F., Campillo, M., Hadziioannou, C., Shapiro, N. M., Nadeau,
R. M., \& Larose, É. (2008a). Postseismic relaxation along the San
Andreas fault at Parkfield. \emph{Science}, 321, 1478--1481.

Brenguier, F., Shapiro, N. M., Campillo, M., Ferrazzini, V., Duputel,
Z., Coutant, O., \& Nercessian, A. (2008b). Towards forecasting volcanic
eruptions using seismic noise. \emph{Nature Geoscience}, 1, 126--130.

Carter, A., van Wyk de Vries, B., Kelfoun, K., Bachèlery, P., \& Briole,
P. (2007). Pits, rifts and slumps: The summit structure of Piton de la
Fournaise. \emph{Bulletin of Volcanology}, 69, 741--756.

Clarke, D., Zaccarelli, L., Shapiro, N. M., \& Brenguier, F. (2011).
Assessment of resolution and accuracy of the moving window cross
spectral technique for monitoring crustal temporal variations using
ambient seismic noise. \emph{Geophysical Journal International}, 186,
867--882.

Clements, T., \& Denolle, M. A. (2018). Tracking groundwater levels
using the ambient seismic field. \emph{Geophysical Research Letters},
45, 6459--6465.

Clements, T., \& Denolle, M. A. (2023). The seismic signature of
California's earthquakes, droughts, and floods. \emph{Journal of
Geophysical Research: Solid Earth}, 128, e2022JB025553.

D'Auria, L., et al.~(2023). Spatio-temporal velocity variations observed
during the pre-eruptive episode of La Palma 2021 eruption inferred from
ambient noise interferometry. \emph{Scientific Reports}, 13, 12203.

Davis, E. E., Sun, T., Heesemann, M., Becker, K., \& Schlesinger, A.
(2024). Long-term offshore borehole fluid-pressure monitoring at the
northern Cascadia subduction zone and inferences regarding the state of
megathrust locking. \emph{Geochemistry, Geophysics, Geosystems}, 24,
e2023GC010910.

Delouche, E., et al.~(2023). Seasonal seismic velocity variations
measured using seismic noise autocorrelations to monitor the dynamic of
aquifers in Greece. \emph{Journal of Geophysical Research: Solid Earth},
128, e2023JB026759.

Donaldson, C., Winder, T., Caudron, C., \& White, R. S. (2019). Crustal
seismic velocity responds to a magmatic intrusion and seasonal loading
in Iceland's Northern Volcanic Zone. \emph{Science Advances}, 5,
eaax6642.

Dvorkin, J., \& Nur, A. (1996). Elasticity of high-porosity sandstones:
Theory for two North Sea data sets. \emph{Geophysics}, 61, 1363--1370.

Ermert, L. A., Cabral-Cano, E., Chaussard, E., et al.~(2023). Probing
environmental and tectonic changes underneath Mexico City with the urban
seismic field. \emph{Solid Earth}, 14, 529--549.

Feng, K.-F., Huang, H.-H., Hsu, Y.-J., \& Wu, Y.-M. (2021). Controls on
seasonal variations of crustal seismic velocity in Taiwan. \emph{Journal
of Geophysical Research: Solid Earth}, 126, e2021JB022650.

Fokker, E., Ruigrok, E., Hawkins, R., \& Trampert, J. (2021).
Physics-based relationship for pore pressure and vertical stress
monitoring using seismic velocity variations. \emph{Remote Sensing}, 13,
2684.

Gassenmeier, M., Sens-Schönfelder, C., Eulenfeld, T., Bartsch, M.,
Victor, P., Tilmann, F., \& Korn, M. (2016). Field observations of
seismic velocity changes caused by shaking-induced damage and healing
due to mesoscopic nonlinearity. \emph{Geophysical Journal
International}, 204, 1490--1502.

Han, S., Bangs, N. L., Carbotte, S. M., Saffer, D. M., \& Gibson, J. C.
(2017). Links between sediment consolidation and Cascadia megathrust
slip behaviour. \emph{Nature Geoscience}, 10, 954--959.

Hotovec-Ellis, A. J., Shiro, B. R., Shelly, D. R., Anderson, K. R.,
Haney, M. M., Thelen, W. A., et al.~(2022). Earthquake-derived seismic
velocity changes during the 2018 caldera collapse of Kīlauea volcano.
\emph{Journal of Geophysical Research: Solid Earth}, 127, e2021JB023324.

Hassanizadeh, S. M., \& Gray, W. G. (1990). Mechanics and thermodynamics
of multiphase flow in porous media including interphase boundaries.
\emph{Advances in Water Resources}, 13, 169--186.

Hassanizadeh, S. M., Celia, M. A., \& Dahle, H. K. (2002). Dynamic
effect in the capillary pressure--saturation relationship. \emph{Vadose
Zone Journal}, 1, 38--57.

Hillers, G., Ben-Zion, Y., Campillo, M., \& Zigone, D. (2015). Seasonal
variations of seismic velocities in the San Jacinto fault area observed
with ambient seismic noise. \emph{Geophysical Journal International},
202, 920--932.

Hobiger, M., Wegler, U., Shiomi, K., \& Nakahara, H. (2016). Coseismic
and post-seismic velocity changes detected by Passive Image
Interferometry: comparison of one great and five strong earthquakes in
Japan. \emph{Geophysical Journal International}, 205, 1053--1073.

Hotovec-Ellis, A. J., et al.~(2022). Earthquake-derived seismic velocity
changes during the 2018 caldera collapse of Kīlauea volcano.
\emph{Journal of Geophysical Research}, 127, e2021JB023324.

Hudson, J. A. (1981). Wave speeds and attenuation of elastic waves in
material containing cracks. \emph{Geophysical Journal International},
64, 133--150.

Hughes, D., \& Kelly, J. L. (1953). Second-order elastic deformation of
solids. \emph{Physical Review}, 92, 1145--1149.

Illien, L., Sens-Schönfelder, C., Andermann, C., Marc, O., Hosseiny, B.,
\& Hovius, N. (2022). Seismic velocity recovery in the subsurface:
Transient damage and groundwater drainage following the 2015 Gorkha
earthquake, Nepal. \emph{Journal of Geophysical Research: Solid Earth},
127, e2021JB023402.

James, S. R., Knox, H. A., Abbott, R. E., \& Screaton, E. J. (2017).
Improved moving window cross-spectral analysis for resolving large
temporal seismic velocity changes in permafrost. \emph{Geophysical
Research Letters}, 44, 4018--4026.

Jeppson, T. N., \& Tobin, H. J. (2015). San Andreas fault zone velocity
structure at SAFOD at core, log, and seismic scales. \emph{Journal of
Geophysical Research: Solid Earth}, 120, 4983--4997.

Johnson, P. A., \& Rasolofosaon, P. (1996). Nonlinear elasticity and
stress-induced anisotropy in rock. \emph{Journal of Geophysical
Research}, 101, 3113--3124.

Kidiwela, M., Denolle, M. A., Wilcock, W. S. D., \& Feng, K.-F. (2026).
Active protothrusts and fluid highways: Seismic noise reveals hidden
subduction dynamics in Cascadia. \emph{Science Advances}, 12, eaea3684.
https://doi.org/10.1126/sciadv.aea3684

Lecocq, T., Longuevergne, L., Pedersen, H. A., Brenguier, F., \&
Stammler, K. (2017). Monitoring ground water storage at mesoscale using
seismic noise. \emph{Scientific Reports}, 7, 14468.

Li, S., Wang, K., Wang, Y., Jiang, Y., \& Dosso, S. E. (2018).
Geodetically inferred locking state of the Cascadia megathrust based on
a viscoelastic Earth model. \emph{Journal of Geophysical Research: Solid
Earth}, 123, 8056--8072.

Lobkis, O. I., \& Weaver, R. L. (2003). Coda-wave interferometry in
finite solids. \emph{Physical Review Letters}, 90, 254302.

Mao, S., Campillo, M., van der Hilst, R. D., Brenguier, F., Stehly, L.,
\& Hillers, G. (2019). High temporal resolution monitoring of small
variations in crustal strain by dense seismic arrays. \emph{Geophysical
Research Letters}, 46, 128--137.

Mao, S., Lecointre, A., van der Hilst, R. D., \& Campillo, M. (2022).
Space-time monitoring of groundwater fluctuations with passive seismic
interferometry. \emph{Nature Communications}, 13, 4643.

Marc, O., Sens-Schönfelder, C., Illien, L., et al.~(2021). Toward using
seismic interferometry to quantify landscape-scale erosion rates.
\emph{Journal of Geophysical Research: Earth Surface}, 126,
e2021JF006112.

Mindlin, R. D. (1949). Compliance of elastic bodies in contact.
\emph{Journal of Applied Mechanics}, 16, 259--268.

Murnaghan, F. D. (1937). Finite deformations of an elastic solid.
\emph{American Journal of Mathematics}, 59(2), 235--260.

Nakata, N., Gualtieri, L., \& Fichtner, A. (2019). \emph{Seismic Ambient
Noise}. Cambridge University Press.

Neal, C. A., \& Lockwood, J. P. (2003). Geologic map of the summit
region of Kilauea Volcano, Hawaii. \emph{U.S. Geological Survey Geologic
Investigations Series}, I-2759.

Niu, F., Silver, P. G., Daley, T. M., Cheng, X., \& Majer, E. L. (2008).
Preseismic velocity changes observed from active source monitoring at
Parkfield. \emph{Nature}, 454, 204--208.

Oakley, D. O. S., et al.~(2021). Seismic ambient noise analyses reveal
changing temperature and water signals to 10s of meters depth in the
critical zone. \emph{Journal of Geophysical Research: Earth Surface},
126, e2020JF005823.

Obermann, A., Planès, T., Larose, E., \& Campillo, M. (2013). Imaging
preeruptive and coeruptive structural and mechanical changes of a
volcano with ambient seismic noise. \emph{Journal of Geophysical
Research}, 118, 6285--6294.

Obermann, A., Planès, T., Larose, E., Sens-Schönfelder, C., \& Campillo,
M. (2014). Depth sensitivity of seismic coda waves to velocity
perturbations in an elastic heterogeneous medium. \emph{Geophysical
Journal International}, 194, 372--382.

Okubo, K., Delbridge, B. G., \& Denolle, M. A. (2024). Monitoring
velocity change over 20 years at Parkfield. \emph{Journal of Geophysical
Research: Solid Earth}, 129, e2023JB028084.

Ostrovsky, L. A., \& Johnson, P. A. (2001). Dynamic nonlinear elasticity
in geomaterials. \emph{La Rivista del Nuovo Cimento}, 24(7), 1--46.

Poupinet, G., Ellsworth, W. L., \& Frechet, J. (1984). Monitoring
velocity variations in the crust using earthquake doublets.
\emph{Journal of Geophysical Research}, 89, 5719--5731.

Richter, T., Sens-Schönfelder, C., Kind, R., \& Asch, G. (2014).
Comprehensive observation and modeling of earthquake and
temperature-related seismic velocity changes in northern Chile.
\emph{Journal of Geophysical Research: Solid Earth}, 119, 4747--4765.

Rivet, D., Campillo, M., Radiguet, M., Dahm, T., \& Shapiro, N. M.
(2011). Seismic evidence of nonlinear crustal deformation during a large
slow slip event in Mexico. \emph{Geophysical Research Letters}, 38,
L08308.

Rodríguez Tribaldos, V., \& Ajo-Franklin, J. B. (2021). Aquifer
monitoring using ambient seismic noise recorded with distributed
acoustic sensing (DAS) deployed on dark fiber. \emph{Journal of
Geophysical Research: Solid Earth}, 126, e2020JB021004.

Roeloffs, E. (1988). Fault stability changes induced beneath a
reservoir. \emph{Journal of Geophysical Research}, 93, 2107--2124.

Sayers, C. M., \& Kachanov, M. (1995). Microcrack-induced elastic wave
anisotropy of brittle rocks. \emph{Journal of Geophysical Research},
100, 4149--4156.

Segall, P., Anderson, K. R., Pulvirenti, F., Wang, T., \& Johanson, I.
(2020). Caldera collapse geometry revealed by near-field GPS
displacements at Kīlauea volcano in 2018. \emph{Geophysical Research
Letters}, 47, e2020GL088867.

Sens-Schönfelder, C., \& Eulenfeld, T. (2019). Probing the in situ
elastic nonlinearity of rocks with Earth tides and seismic noise.
\emph{Physical Review Letters}, 122, 138501.

Sens-Schönfelder, C., \& Wegler, U. (2006). Passive image interferometry
and seasonal variations of seismic velocities at Merapi Volcano,
Indonesia. \emph{Geophysical Research Letters}, 33, L21302.

Shapiro, S. A. (2003). Elastic piezosensitivity of porous and fractured
rocks. \emph{Geophysics}, 68, 482--486.

Shi, Q., Montgomery, D. R., Swann, A. L. S., Cristea, N. C., Williams,
E. F., You, N., Jeffery, S., Collins, J., Prada Barrio, A., Misiewicz,
P. A., Nissen-Meyer, T., \& Denolle, M. A. (2026). Agroseismology and
the impact of farming practices on soil hydrodynamics. \emph{Science},
392, 306--310. https://doi.org/10.1126/science.aec0970

Singh, J., Curtis, A., Zhao, Y., Cartwright-Taylor, A., \& Main, I.
(2019). Coda wave interferometry for accurate simultaneous monitoring of
velocity and acoustic source location changes. \emph{Journal of
Geophysical Research: Solid Earth}, 124, 5629--5655.

Smail, T., et al.~(2019). Cumulative deviation of precipitation for
estimating groundwater levels. \emph{Water Resources Research}, 55,
10672--10689.

Snieder, R. (2002). Coda wave interferometry and the equilibration of
energy in elastic media. \emph{Physical Review E}, 66, 046615.

Snieder, R. (2006). The theory of coda wave interferometry. \emph{Pure
and Applied Geophysics}, 163, 455--473.

Snieder, R., Grêt, A., Douma, H., \& Scales, J. (2002). Coda wave
interferometry for estimating nonlinear behavior. \emph{Science}, 295,
2253--2255.

Snieder, R., Sens-Schönfelder, C., Ruigrok, E., \& Shiber, K. (2017).
Seismic velocity changes in the shallow subsurface due to the non-linear
compliance of fractured rock. \emph{Geophysical Journal International},
211, 1487--1494.

Taira, T., Nayak, A., Brenguier, F., \& Manga, M. (2015). Monitoring
reservoir response to earthquakes and fluid extraction. \emph{Science
Advances}, 4, e1701536.

Takano, T., Nishimura, T., Nakahara, H., Ueda, H., \& Fujita, E. (2019).
Sensitivity of seismic velocity changes to the tidal strain at different
lapse times: Data analyses of a small seismic array at Izu-Oshima
volcano. \emph{Journal of Geophysical Research: Solid Earth}, 124,
3011--3023. https://doi.org/10.1029/2018JB016235

Talwani, P., Chen, L., \& Gahalaut, K. (2007). Seismogenic permeability.
\emph{Journal of Geophysical Research}, 112, B07309.

Tromp, J., \& Trampert, J. (2018). Effects of induced stress on seismic
forward modelling and inversion. \emph{Geophysical Journal
International}, 213, 851--867.

Tromp, J., Tape, C., \& Liu, Q. (2005). Seismic tomography, adjoint
methods, time reversal and banana-doughnut kernels. \emph{Geophysical
Journal International}, 160, 195--216.

Tsai, V. C. (2011). A model for seasonal changes in GPS positions and
seismic wave speeds. \emph{Journal of Geophysical Research}, 116,
B04404.

Verdon, J. P., Angus, D. A., Kendall, J. M., \& Hall, S. A. (2008). The
effect of microstructure and nonlinear stress on anisotropic seismic
velocities. \emph{Geophysics}, 73, D41--D51.

Walsh, J. B. (1965). The effect of cracks on the compressibility of
rock. \emph{Journal of Geophysical Research}, 70, 381--389.

Wegler, U., \& Sens-Schönfelder, C. (2007). Fault zone monitoring with
passive image interferometry. \emph{Geophysical Journal International},
168, 1029--1033.

Yuan, C., Bryan, J., \& Denolle, M. A. (2021). Numerical comparison of
time-, frequency- and wavelet-domain methods for coda wave
interferometry. \emph{Geophysical Journal International}, 226, 828--846.

Zhan, Z., Tsai, V. C., \& Clayton, R. W. (2013). Spurious velocity
changes caused by temporal variations in ambient noise frequency
content. \emph{Geophysical Journal International}, 194, 1552--1559.

Zhang, S., Luo, B., Ben-Zion, Y., Lumley, D. E., \& Zhu, H. (2023).
Monitoring terrestrial water storage, drought and seasonal changes in
central Oklahoma with ambient seismic noise. \emph{Geophysical Research
Letters}, 50, e2023GL103419.

Zhang, H., Thurber, C., \& Bedrosian, P. (2009). Joint inversion for Vp,
Vs, and Vp/Vs at SAFOD, Parkfield, California. \emph{Geochemistry,
Geophysics, Geosystems}, 10, Q11002.

Zhu, T., Ajo-Franklin, J., Daley, T. M., \& Marone, C. (2019). Dynamics
of geologic CO\(_2\) storage and plume motion revealed by seismic coda
waves. \emph{Proceedings of the National Academy of Sciences}, 116,
2464--2469.

\clearpage

\section{Table Captions}\label{table-captions}

\textbf{Table 1.} Key parameters and validity limits for applying the
framework. The full parameter overview, including additional ranges and
citations, is provided in Table S1.

\clearpage

\section*{Figure Files}

\clearpage
\begin{figure}[p]
\centering
\includegraphics[width=\textwidth,height=0.72\textheight,keepaspectratio]{figures/main/fig01_unified_workflow.png}
\caption{Unified workflow for interpreting ambient-noise velocity changes. The workflow connects external forcing (temperature, hydrology, surface loading, tectonics, and magmatic pressure) to stress and strain perturbations, material sensitivity parameters ($\beta$, $\mu'$, crack fabric, saturation state, and rheology), frequency-dependent $\delta v/v(f,t)$, and the diagnostic steps targeted for implementation in \texttt{codameter}.}
\end{figure}

\clearpage
\begin{figure}[p]
\centering
\includegraphics[width=\textwidth,height=0.72\textheight,keepaspectratio]{figures/main/fig02_depth_kernels.png}
\caption{Frequency-dependent depth sensitivity and inversion limits. (a) Schematic Rayleigh/Love sensitivity kernels showing broad, overlapping depth sensitivity for representative frequency bands. (b) Peak sensitivity depth $z_{\mathrm{peak}} \approx V_S/(3f)$ for soft sediment, weathered rock, and crystalline rock. (c) Relative uncertainty reduction for independent versus correlated frequency bands, emphasizing that depth-resolved inversion requires regularization, uncertainty propagation, and resolution tests.}
\end{figure}

\clearpage
\begin{figure}[p]
\centering
\includegraphics[width=\textwidth,height=0.72\textheight,keepaspectratio]{figures/main/fig03_hydrological_competition.png}
\caption{Hydrological signals are a sign competition between surface loading and pore pressure. The simplified Fokker et al. (2021) endmember shows the opposing contributions for (a) ice/snow loading, (b) seasonal rainfall, and (c) reservoir impoundment. Positive-compressive vertical loading tends to increase velocity, whereas pore-pressure increase reduces effective stress and decreases velocity.}
\end{figure}

\clearpage
\begin{figure}[p]
\centering
\includegraphics[width=\textwidth,height=0.72\textheight,keepaspectratio]{figures/main/fig04_material_sensitivity.png}
\caption{Material sensitivity controls stress conversion from velocity change. (a) Effective $|\beta|$ spans orders of magnitude across materials, from intact crystalline rock to soft sediment. (b) Observed $\delta v/v$ amplitude ratios across Parkfield, Cascadia, and Kīlauea do not map directly to stress ratios because the conversion depends on $\beta$, $\mu'$, and shear modulus.}
\end{figure}

\clearpage
\begin{figure}[p]
\centering
\includegraphics[width=\textwidth,height=0.72\textheight,keepaspectratio]{figures/main/fig05_anisotropy_fabric.png}
\caption{Deviatoric stress and fracture fabric determine the observed velocity sign. (a) Directional sensitivity increases with fracture fabric strength. (b–d) Parkfield, Cascadia, and Kīlauea select different stress/strain components because the dominant loading geometry and crack fabric differ. (e) The scalar isotropic model succeeds for Cascadia but fails for Parkfield and Kīlauea.}
\end{figure}

\clearpage
\begin{figure}[p]
\centering
\includegraphics[width=\textwidth,height=0.72\textheight,keepaspectratio]{figures/main/fig06_rheology_diagnostics.png}
\caption{Rheology is diagnosed from phase, hysteresis, and recovery shape. (a) Endmember elastic, Maxwell, Kelvin-Voigt, and slow-dynamics responses to a step-like forcing. (b) $\delta v/v$–strain crossplots distinguish elastic, higher-order nonlinear, and viscoelastic lag behavior. (c) Parkfield-like dual-population model combining tectonic loading and postseismic healing. (d) Tidal probing of in-situ nonlinearity through harmonic distortion.}
\end{figure}

\end{document}
